\section{Averaged holonomy and the category argument}
\label{ho:section}

We turn the finite curvature difference of
\Cref{prop:curvature-signal} into an obstruction to a square-integrable
connection.  The estimate is averaged over four position variables and
three momentum variables.  The loops themselves remain in the smooth
development.  In particular, no restriction of the weak ambient
connection to a prescribed curve is required.
Sbierski developed holonomy obstructions at Lipschitz regularity
\cite{Sbierski2022} and a related nonsymmetric obstruction based on
integrated curvature \cite{Sbierski2026}. Here the ambient hypothesis controls
an integral over positions; averaging must therefore precede any
curvewise transport estimate. The stopping argument below establishes
the needed chart control for the sampled loops.

\subsection{Canonical parameters and the observation frame}

Write
\begin{equation}\label{ho:scales}
 r=r_T=e^{-\kappa T},\qquad l_T=e^{-\kappa T/16},
 \qquad \amp=e^{-\kappa T/4}.
\end{equation}
The observation coordinates are $(u,v,\theta^1,\theta^2)$, and the frame
used throughout this section is
\begin{equation}\label{ho:frame}
 F_0=\partial_u,\qquad F_1=r^{-1}\partial_v,
 \qquad F_2=\partial_{\theta^1},\quad F_3=\partial_{\theta^2}.
\end{equation}
For each fixed experiment, $r$ is a constant on the entire observation
box.  Thus the four fields in \eqref{ho:frame} commute.

Let $O_0\Subset O$ be the closed trimmed ball from
\Cref{prop:matched-seeds}, chosen so that
\begin{equation}\label{ho:trim}
                   \mu_d(O_0)>.95\mu_d(O).
\end{equation}
For $v\in T+J_0$, where $J_0$ is a fixed compact interval of positive
length, let $\mathcal P_{T,v}$ be the set of terminal states
$(u,\theta,\pi_u,P)$ obtained by following the background seeds in $O_0$
to that level.  The full parameter set and its measure are
\begin{equation}\label{ho:parameter-measure}
 \mathcal P_T=\{(v,u,\theta,\pi_u,P):v\in T+J_0,
                \ (u,\theta,\pi_u,P)\in\mathcal P_{T,v}\},\qquad
 \dd\nu_T=\dd v\,\dd u\,\dd^2\theta\,\dd\pi_u\,\dd^2P.
\end{equation}
The three momentum labels range in a fixed bounded box $Q\subset\R^3$.
Here and below, measures may be completed without changing any estimate.

We recall why these are canonical measures, also for the nearby metrics
used later.  On the future unit shell of
$\mathcal H(x,\pi)=\frac12g^{\alpha\beta}(x)\pi_\alpha\pi_\beta$,
a transverse $v$ section has a simple root for $\pi_v$ because
$\partial_{\pi_v}\mathcal H=\dot v>0$.  On that section the pullback of
the symplectic form is
\[
 \dd\pi_u\wedge\dd u+\sum_{A=1}^2\dd P_A\wedge\dd\theta^A.
\]
The Hamiltonian vector field is the kernel of the symplectic form
restricted to the shell.  Its transverse hitting maps consequently
preserve this two-form and its canonical volume, the absolute value of
its third exterior power divided by $3!$.  The corresponding
section on $\Sigma$, pulled back by $w\mapsto h(w,\cdot)$, is precisely
$\mu_d$.  Hence
\begin{equation}\label{ho:slice-volume}
 \int_{\mathcal P_{T,v}}\dd u\,\dd^2\theta\,\dd\pi_u\,\dd^2P
       =\mu_d(O_0),\qquad
 \nu_T(\mathcal P_T)=|J_0|\mu_d(O_0).
\end{equation}
This uses injectivity of transverse phase-space hitting maps; it makes
no injectivity assertion about their projection onto positions.

The following argument applies separately to the background metric and
to any persisted finite experiment satisfying
\Cref{prop:matched-seeds}.  Denote either metric temporarily by $g$ and
put $G_{ij}=g(F_i,F_j)$.  On an observation box containing the positions
in \eqref{ho:parameter-measure} with a fixed margin, we have
\begin{equation}\label{ho:gram-assumptions}
 \|G\|+\|G^{-1}\|\le C,
 \qquad
 \sum_{|\beta|\le30}
 \bigl(\|\partial^\beta G\|+\|\partial^\beta G^{-1}\|\bigr)
       \le Q_T,\qquad Q_T=e^{o(T)}.
\end{equation}
All derivatives here are ordinary coordinate derivatives.
The selected future unit tangent $U$ at every parameter in
$\mathcal P_T$ has bounded $F$ components.  Constants in this section
may depend on the fixed observation tube, $O$, and the chart bound $B$,
but not on $T$.  Finite products of subexponential bounds are absorbed
into a new $Q_T$ as in \Cref{conv:bg-subexp}.

If $g_{\rm coord}$ is the ordinary coordinate matrix, the constant
frame change gives
$\det G=r^{-2}\det g_{\rm coord}$.  Thus
\begin{equation}\label{ho:volume-density}
 \dd\operatorname{vol}_g
     =r\sqrt{|\det G|}\,\dd u\,\dd v\,\dd^2\theta,
 \qquad c r\le\sqrt{|\det g_{\rm coord}|}\le C r.
\end{equation}
For an optical metric these bounds follow directly from
$a/r\asymp1$, the bounded angular metric and shift.  Formula
\eqref{ho:volume-density} also holds for the nearby metrics without
requiring the observation coordinates to remain optical.

\begin{lemma}[Comparison at a good base point]\label{ho:base-comparison}
Suppose an extension chart has bounded metric and inverse metric, and
the components of $U$ are bounded both in that chart and in $F$.
Then the matrix $M$ with columns $F_i$ in chart coordinates and its
inverse are bounded by a constant depending only on these bounds and
\eqref{ho:gram-assumptions}.
\end{lemma}

\begin{proof}
For any Lorentzian Gram matrix $G$ and unit timelike column $U$, set
\[
                 K(G,U)=G+2(GU)(GU)^{\mathsf t}.
\]
The associated form is positive definite: writing a vector as
$aU+Z$ with $g(U,Z)=0$ gives $a^2+g(Z,Z)>0$ unless the vector vanishes.
The determinant lemma and $U^{\mathsf t}GU=-1$ give the exact identity
\begin{equation}\label{ho:positive-determinant}
       \det K(G,U)=\det G\,(1+2U^{\mathsf t}GU)=-\det G.
\end{equation}
Bounded $G,G^{-1},U$ give an upper eigenvalue bound for $K$ and a
positive lower bound for its determinant.  Its smallest eigenvalue is
therefore bounded below by the determinant divided by the cube of the
upper eigenvalue bound.  Apply this in both bases.  If $g_c$ is the
chart matrix and $U_c=MU_F$, then
\[
 K(G,U_F)=M^{\mathsf t}K(g_c,U_c)M.
\]
The two uniform positive definite comparisons bound both $M$ and
$M^{-1}$.  In particular, bounded Lorentzian Gram matrices alone are
not being used to bound a Lorentz boost.
\end{proof}

\subsection{Finite jets of coordinate rectangles}

The scale choices can be read from the estimates proved below. With a
stopped edge-integral cutoff $R$, position averaging will leave
exceptional measure $O((rR^2)^{-1})$, while the extracted quadratic
coefficient will be bounded by
$O(rR/l_T)+Q_Tl_T^{16}$. Taking
\[
 R=r^{-17/32},\qquad l_T=r^{1/16},\qquad \amp=r^{1/4}
\]
gives exceptional measure $O(r^{1/16})$ and curvature errors
$O(r^{13/32})+Q_Tr=o(\amp)$. The power $l_T^{16}$ comes from an
eighteenth-order remainder after extracting a quadratic coefficient.

Let
\begin{equation}\label{ho:directions}
 \mathcal X=\{\partial_{\theta^1},\partial_{\theta^2},
                   \partial_{\theta^1}+\partial_{\theta^2}\}.
\end{equation}
For a position $x$ in the parameter set, $X\in\mathcal X$, and
$0\le l\le17l_T$, take the coordinate rectangle with consecutive
displacements $l\partial_v,lrX,-l\partial_v,-lrX$.
Denote its holonomy in the based frame $F(x)$ by $H_{x,X}(l)$.
Every rectangle lies in the enlarged observation box for large $T$,
because its ordinary coordinate displacement is $O(l)$.
Its velocity in $F$, on each unit-parameter edge, is $lr$ times a fixed
bounded column.  The two ordinary side lengths have different scales,
although both directions have this same $F$-component scale.
\Cref{ho:fig-loops} illustrates these two scales.
The quadratic coefficient of the loop holonomy recovers the curvature
block $D$. We will extract that coefficient from the eighteen values
at $l=0,l_T,\ldots,17l_T$; a bound for the eighteenth derivative
controls the interpolation remainder at the required $o(\amp)$ scale.
\begin{figure}[htbp]
\centering
\begin{tikzpicture}[>=Latex,font=\small,line cap=round]
 \begin{scope}[shift={(0,0)}]
  \node[font=\small\bfseries] at (2,2.65) {Coordinate displacements};
  \draw[->,gray] (-.25,0)--(4.65,0) node[right] {$\Delta v$};
  \draw[->,gray] (0,-.25)--(0,1.9) node[above] {$\Delta\theta^1$};
  \draw[blue!65!black,thick,->] (0,0)--(2,0);
  \draw[blue!65!black,thick] (2,0)--(4,0);
  \draw[blue!65!black,thick,->] (4,0)--(4,.8);
  \draw[blue!65!black,thick] (4,.8)--(4,1.3);
  \draw[blue!65!black,thick,->] (4,1.3)--(2,1.3);
  \draw[blue!65!black,thick] (2,1.3)--(0,1.3);
  \draw[blue!65!black,thick,->] (0,1.3)--(0,.55);
  \draw[blue!65!black,thick] (0,.55)--(0,0);
  \fill (0,0) circle (1.8pt) node[below left] {$x$};
  \node[below] at (2,0) {$l$};
  \node[right] at (4,.65) {$rl$};
 \end{scope}
 \begin{scope}[shift={(7,0)}]
  \node[font=\small\bfseries] at (1.25,2.65) {$F$ components};
  \draw[->,gray] (-.25,0)--(3.25,0) node[right] {$r\Delta v$};
  \draw[->,gray] (0,-.25)--(0,1.9) node[above] {$\Delta\theta^1$};
  \draw[blue!65!black,thick,->] (0,0)--(.9,0);
  \draw[blue!65!black,thick] (.9,0)--(1.5,0);
  \draw[blue!65!black,thick,->] (1.5,0)--(1.5,.9);
  \draw[blue!65!black,thick] (1.5,.9)--(1.5,1.5);
  \draw[blue!65!black,thick,->] (1.5,1.5)--(.6,1.5);
  \draw[blue!65!black,thick] (.6,1.5)--(0,1.5);
  \draw[blue!65!black,thick,->] (0,1.5)--(0,.6);
  \draw[blue!65!black,thick] (0,.6)--(0,0);
  \fill (0,0) circle (1.8pt);
  \node[below] at (.75,0) {$rl$};
  \node[right] at (1.5,.75) {$rl$};
 \end{scope}
\end{tikzpicture}
\caption{A coordinate rectangle for $X=\partial_{\theta^1}$, with $u$
and $\theta^2$ fixed (schematic).  Its sides have equal component scale
in $F$, even though their coordinate scales differ.  The proof averages
translated rectangles over all four base-position variables and the
three bounded momentum labels.}
\label{ho:fig-loops}
\end{figure}


\begin{lemma}[Loop jets and the quadratic coefficient]\label{ho:loop-jets}
Uniformly in the observation positions and $X\in\mathcal X$,
\begin{equation}\label{ho:jet-bound}
 \sup_{0\le l\le17l_T}\|\partial_l^{18}H_{x,X}(l)\|\le Q_T.
\end{equation}
The quadratic Taylor coefficient $C_2(x,X)$ of $H_{x,X}(l)-\Id$ is
the matrix, in $F(x)$, of $\pm R(\partial_v,rX)$.  Consequently
\begin{equation}\label{ho:curvature-cancellation}
 g\bigl(C_2(x,X)F_1,X\bigr)
       =\pm R(\partial_v,X,\partial_v,X)=\pm D(X,X).
\end{equation}
The same curvature and loop-orientation convention fixes the signs.
\end{lemma}

\begin{proof}
Since $F$ commutes, its connection coefficients are given by Koszul's
formula
\[
 C^k{}_{ij}=\tfrac12G^{kh}
       (F_iG_{jh}+F_jG_{ih}-F_hG_{ij}).
\]
For a constant column $V$ and $W=rV^iF_i$, the connection matrix on
$W$ is therefore
\begin{equation}\label{ho:koszul-scaled}
 \mathsf A(W)^k{}_j
 =\tfrac12G^{kh}V^i
       \bigl((rF_i)G_{jh}+(rF_j)G_{ih}-(rF_h)G_{ij}\bigr).
\end{equation}
Each $rF_i$ is one of $r\partial_u,\partial_v,r\partial_{\theta^A}$.
Their coefficients are constant and bounded.  Thus the connection
contracted with each of $\partial_v$ and $rX$, and its ordinary
derivatives through order eighteen, are bounded by $Q_T$.
This proves the required cancellation directly, including all shift
and lapse terms; uncontracted $F$ connection coefficients need not
have such a bound.

For an edge parametrized by $s\in[0,1]$, its position has the form
$x+c_e(s,l)$, where $c_e$ is linear in $l$ and independent of momentum.
Its transport equation in $F$ has coefficient
$l\mathsf A(\pm W)(x+c_e(s,l))$ for $W=\partial_v$ or $rX$.
Its integral norm is at most $lQ_T=o(1)$, uniformly in the indicated
range.  Differentiating this coefficient at most eighteen times in
$l$ uses only ordinary derivatives in bounded coordinate directions,
so each derivative is bounded by $Q_T$.
To see the corresponding assertion for transport, differentiate
$P'=-A_eP$, $P(0)=\Id$, $k$ times in $l$:
\[
 (\partial_l^kP)'=-A_e\partial_l^kP
       -\sum_{j=1}^k\binom{k}{j}
            (\partial_l^jA_e)(\partial_l^{k-j}P).
\]
Gronwall uses only the small integral of the undifferentiated $A_e$.
Induction in $k\le18$ bounds every differentiated solution by finite
products of $Q_T$, hence by another subexponential bound.  Multiplying
the four edge transports proves \eqref{ho:jet-bound}.

For completeness, pull the smooth connection back by
$(s,t)\mapsto x+s\partial_v+trX$.  Choose a smooth frame agreeing
with $F$ at $(0,0)$ whose connection one-form vanishes there.
This is possible by prescribing the first derivatives of a frame
change at that point.  On the boundary of $[0,l]^2$ its connection
one-form has integral norm $O(l^2)$.  The ordered transport equals
$\Id$ minus its boundary integral up to $O(l^4)$.  Stokes' theorem
identifies the quadratic coefficient with minus its exterior
derivative at the origin; since the connection form vanishes there,
this is the curvature on $(\partial_v,rX)$, with the orientation sign.
This auxiliary frame is used only to identify a Taylor coefficient,
not to obtain a uniform estimate.
Finally, multilinearity gives
\[
 g\bigl(R(\partial_v,rX)(r^{-1}\partial_v),X\bigr)
       =r r^{-1}R(\partial_v,X,\partial_v,X).
\]
This proves \eqref{ho:curvature-cancellation} with no remaining
factor of $r$.
\end{proof}

\subsection{Averaging before closing the chart estimates}

\begin{proposition}[Averaged holonomy obstruction]
\label{prop:averaged-holonomy}
Assume \eqref{ho:gram-assumptions} on the observation box, bounded
$F$ components of the selected unit tangents, and a fixed bounded
momentum box in \eqref{ho:parameter-measure}. Fix one extension and
one collection of at most $B$ charts with the metric, inverse-metric,
and connection bounds in \Cref{def:exit-test}.
Let $\mathcal G_T\subset\mathcal P_T$ be the measurable parameters
whose base point and unit tangent satisfy the margin and velocity
conditions in at least one chart of this same collection.
There is a measurable exceptional set $\mathcal N_T\subset\mathcal G_T$
such that
\begin{equation}\label{ho:exception-bound}
 \nu_T(\mathcal N_T)\le C e^{-\kappa T/16},
\end{equation}
and, for all parameters in $\mathcal G_T\setminus\mathcal N_T$,
\begin{equation}\label{ho:D-small}
 |D|\le C e^{-13\kappa T/32}+Q_Te^{-\kappa T}
                 =o(\amp).
\end{equation}
The constants use only $B$ and the uniform observation bounds.
\end{proposition}

\begin{proof}
At each good parameter select the least chart index for which the
coordinate margin is at least $1/B$ and $|U_c|\le B$.
The finite selector is measurable and may depend on all three momentum
labels. Fix Borel representatives of the finitely many weak Christoffel
arrays. By \Cref{ho:base-comparison}, the base frame matrices obey
$\|M(0)\|+\|M(0)^{-1}\|\le C_0$.
Choose a fixed $K>2C_0+1$.

Fix $X\in\mathcal X$ and a sampled positive size
$l=j l_T$, $1\le j\le17$.  Use the selected base chart throughout the
four concatenated edges.  Define one first stop on this concatenation:
stop when either $\|M\|=K$ or the chart position has moved Euclidean
distance $1/(2B)$ from its base position.  The latter displacement
keeps the path in the chart because the base margin is $1/B$.
An edge after an earlier stop contributes nothing to any integral
below. Before this global stop the smooth interior path, frame matrix,
and chart coordinates are all defined. More precisely, a parameter is
before the stop at a given time when both strict thresholds hold at
every earlier time in the concatenation. In a fixed chart, continuity
expresses this condition using suprema over rational earlier times, so
it is Borel on the smooth state domain. Intersecting with the measurable
selected parameter set gives the stopped subsets used below.

Write $x=(v,u,\theta)$ and $m=(\pi_u,P)$.  At a fixed edge parameter
$s$ the position is the translation $x\mapsto x+c_e(s,l)$, independent
of $m$.  For the four edges the translations are, successively,
\[
 sl\partial_v,\quad l\partial_v+slrX,\quad
 (1-s)l\partial_v+lrX,\quad (1-s)lrX.
\]
For a selected chart $c$, restrict first to good parameters before
the global stop.  Enlarge this nonnegative integral to all $m\in Q$
and all translated observation positions lying in that chart.  On
positions the translation, the smooth open isometric embedding, and
the chart map compose to an injective smooth coordinate map.  Its
volume change is controlled without estimating its derivative matrix:
isometry and \eqref{ho:volume-density} give
\begin{equation}\label{ho:density-only}
 \dd x\le C r^{-1}\dd\operatorname{vol}_g
       =C r^{-1}\sqrt{|\det\widetilde g_c|}\,\dd y_c
       \le C_B r^{-1}\dd y_c.
\end{equation}
Consequently, for each fixed $s$,
\begin{equation}\label{ho:fixed-edge-average}
 \int_{\substack{\mathcal G_T:\text{chart }c\text{ selected}\\
                         \text{before global stop}}}
 |\widetilde\Gamma_c(y_c(x+c_e(s,l)))|^2\,\dd\nu_T
 \le \frac{C_B|Q|}{r}
                  \int_{c}|\widetilde\Gamma_c(y)|^2\,\dd y.
\end{equation}
Here $\dd\nu_T=\dd x\,\dd m$; no seed Jacobian remains.
The bounded momentum fiber contributes only $|Q|$.
Momentum-dependent chart selection and stop conditions have only
restricted the integral before it was enlarged.  In particular,
\eqref{ho:density-only} has not used the unproved along-edge bound for
$M$.

The weak chart Christoffel array agrees almost everywhere on the open
image with its smooth interior array.  For each $s$, the inverse image
of the exceptional position null set under the preceding smooth
coordinate map is null.  Fubini in $s,x,m$ shows that, outside a null
set of parameters, the stopped trajectory integrals are unchanged by
this choice of representative.  We discard that null set for each of
the finitely many sampled loops.  This justifies the use of the
smooth connection equation below, without asserting a Sobolev trace
on an individually prescribed loop.

Integrate \eqref{ho:fixed-edge-average} in $s\in[0,1]$ and sum the
at most $B$ charts.  If
\[
 I_e(x,m)=\int_0^1\boldsymbol 1_{\{\text{before global stop}\}}
                 |\widetilde\Gamma_c(y_e(s))|\,\dd s,
\]
Cauchy--Schwarz and the chart $L^2$ bounds imply
\begin{equation}\label{ho:L2-average}
             \int_{\mathcal G_T} I_e^2\,\dd\nu_T\le C/r.
\end{equation}
Set $R_T=r^{-1/2}e^{\kappa T/32}$.
Markov's inequality applied to $I_e^2$ gives
\begin{equation}\label{ho:markov}
 \nu_T\{I_e>R_T\}
       \le (C/r)R_T^{-2}=C e^{-\kappa T/16}.
\end{equation}
There are four edges, seventeen positive sizes, and three choices of
$X$.  The union of their exceptional sets, together with the null
sets already removed, satisfies \eqref{ho:exception-bound}.
Only this finite collection is used: no common exceptional set for a
continuum of positive loop sizes is asserted.

We now prove that no stop occurs for a remaining parameter.
On one edge write its $F$-velocity as $f=lrV$, with $V$ a fixed bounded
column, and let $\mathsf A$ be its internal connection matrix.
By \eqref{ho:koszul-scaled}, $|\mathsf A|\le lQ_T$.
The connection transformation law along the smooth interior curve is
the exact matrix equation
\begin{equation}\label{ho:frame-ode}
            M'=M\mathsf A-\widetilde\Gamma_c(Mf)M.
\end{equation}
Here $\widetilde\Gamma_c(Mf)$ means the chart connection matrix
contracted with the chart velocity $Mf$.
Before the stop,
\[
 |M'|\le K lQ_T+C K^2lr|\widetilde\Gamma_c|,
                  \qquad |y'|\le C K lr.
\]
Summing over all four stopped edges and using $I_e\le R_T$ yields
\begin{align}\label{ho:stop-improvement}
 \sup\|M-M(0)\|&\le 4K lQ_T+C K^2lrR_T=o(1),\notag\\
 \sup|y-y(0)|&\le C K lr=o(1).
\end{align}
Indeed $l\le17e^{-\kappa T/16}$ and
$lrR_T\le17e^{-17\kappa T/32}$.
For large $T$ these inequalities are strict improvements of both stop
thresholds.  Continuity rules out a first stop, including on any of
the later three edges.  Thus the entire loop stays in its selected
base chart, with $\|M\|<K$.  Only an upper matrix threshold was needed;
the inverse at the base is already bounded.  If desired, inverse
bounds along the loop also follow from
$\det(M)^2=\det G/\det\widetilde g_c$ and the cofactor formula.

Ordinary chart parallel transport now has total integrated coefficient
at most $C lr\sum_e I_e\le C lr^{1/2}e^{\kappa T/32}$.
Its matrix differs from $\Id$ by at most the exponential of this
quantity minus one.  Since the rectangle closes exactly, its final
frame matrix equals its base matrix.  Thus if $P_c$ is the full chart
transport, $H=M(0)^{-1}P_cM(0)$.  Conjugating by the bounded base matrix
and its inverse gives
\begin{equation}\label{ho:sampled-holonomy}
 \|H_{x,X}(j l_T)-\Id\|
       \le C r^{1/2}l_T e^{\kappa T/32},
                     \qquad 1\le j\le17.
\end{equation}
At $j=0$ the left side is zero.  All transports in this argument have
taken place in the smooth original spacetime.

It remains to extract the quadratic coefficient from these eighteen
values.  We include the elementary interpolation argument, as in
\cite[\CompanionRef{signed}{holo:interpolation}]{KerrCOne}.  Write
\[
 H_{x,X}(l)-\Id=\sum_{k=0}^{17}C_k l^k+\mathcal R_{18}(l).
\]
By \Cref{ho:loop-jets}, the Taylor integral remainder at $l=j l_T$
is bounded by $C Q_Tl_T^{18}$.  The Vandermonde matrix
$(j^k)_{0\le j,k\le17}$ is invertible and independent of $T$.
Therefore its inverse expresses $l_T^2 C_2$ as a fixed linear
combination of the values of the polynomial at these nodes, giving
\begin{align}\label{ho:coefficient-estimate}
 \|C_2\|
 &\le C l_T^{-2}\max_{0\le j\le17}
                       \|H_{x,X}(j l_T)-\Id\|+Q_Tl_T^{16}\notag\\
 &\le C r^{1/2}l_T^{-1}e^{\kappa T/32}+Q_Tl_T^{16}\notag\\
 &=C e^{-13\kappa T/32}+Q_Te^{-\kappa T}.
\end{align}
Both exponents have a strict gain over $-\kappa/4$.
By \eqref{ho:curvature-cancellation} and boundedness of $G$, this
controls $D(X,X)$ for each $X\in\mathcal X$.
Curvature pair symmetry gives $D_{12}=D_{21}$, and
\[
 2D_{12}=D(\partial_{\theta^1}+\partial_{\theta^2},
                \partial_{\theta^1}+\partial_{\theta^2})
             -D_{11}-D_{22}.
\]
These three estimates prove \eqref{ho:D-small} for the entire
symmetric $2\times2$ block.
\end{proof}

\subsection{Finite persistence and nowhere denseness}

The preceding estimate makes the observed curvature block small on
almost every parameter with an extension-chart witness. To contradict
the packet signal, we will apply it to the original datum and a nearby
datum at identical numerical terminal states. Canonical section volume
ensures that their two good families still overlap.

\Cref{prop:matched-seeds} supplies the needed persistence for each
fixed finite experiment: a relative open neighborhood preserves the
signal with lower bound $\amp$, the frame and finite-jet bounds, and
the injective measure-preserving backward maps into $O$. Its map
tolerance can be chosen after $T$, and its data-neighborhood radius
may depend on that experiment. The next argument uses precisely this
finite persistence.

\begin{theorem}[Nowhere denseness of the extension tests]
\label{thm:nowhere-dense}
For the single neighborhood $\Ueps$ fixed in
\Cref{prop:bg-package}, every set $E(O,B)$ of \Cref{def:exit-test} is
relatively nowhere dense in $\Ueps$.
\end{theorem}

\begin{proof}
Suppose instead that a nonempty relative open set $V\subset\Ueps$
is contained in $\overline{E(O,B)}$.  Since $E(O,B)$ is dense in $V$,
choose $d\in E(O,B)\cap V$.
The weighted seminorms are increasing, so there are a finite $m\ge10$
and $\eta>0$ such that
\begin{equation}\label{ho:finite-neighborhood}
        \{d'\in\Ueps:p_m(d'-d)<\eta\}\subset V.
\end{equation}
Fix the branch tube and the extension witness belonging to this datum.
Apply the finite packet construction at order $m$, along its unbounded
sequence of admissible observation times.  By
\Cref{prop:packet-preparation,thm:dn-comparison,prop:experiment-attachment},
the exact full-bridge data $d_T$ satisfy $p_m(d_T-d)\to0$; hence they
lie in $V$ for all sufficiently large $T$.

The set $E(O,B)$ is dense in $V$, so \Cref{prop:matched-seeds}
allows us, for every such fixed $T$, to choose
$\widehat d_T\in E(O,B)\cap V$ as close to $d_T$ as required by all the
finite observation conditions.  We also require
$p_0(\widehat d_T-d_T)\to0$.  Write $\widehat g_T$ for its observed
metric.  We then have, on the common position tube,
\begin{equation}\label{ho:persisted-signal}
 |D(\widehat g_T)-D(g_d)|>\amp
                 \quad\text{if }|\cos(\lambda v)|\ge1/2,
\end{equation}
and both metrics have the uniform hypotheses of
\Cref{prop:averaged-holonomy}.  In particular the constants there are
independent of $T$; the required finite closeness was chosen to retain
them.  Smooth dependence of the initial section density on $h$ gives
\begin{equation}\label{ho:initial-measure-ratio}
                  \frac{\mu_{\widehat d_T}(O)}{\mu_d(O)}\longrightarrow1.
\end{equation}

Let $S\subset O$ and $\widehat S_T\subset O$ be the good seed families
in the two existential extension witnesses.  Each has more than
$99$ percent of its respective measure of $O$.
Use \eqref{ho:parameter-measure} for the background hits of $O_0$.
At a fixed $v$, the loss from parameters whose background seed is
outside $S$ is at most $.01\mu_d(O)$.
The matched backward map for $\widehat g_T$ is injective into $O$ and
preserves its canonical section measure.  Its loss from
$O\setminus\widehat S_T$ is therefore at most
$.01\mu_{\widehat d_T}(O)$.
By \eqref{ho:initial-measure-ratio}, their combined loss is less than
$.04\mu_d(O)$ for large $T$, uniformly on every $v$ slice.  Thus, if
$\mathcal G_T^{\cap}$ denotes parameters good for both metrics,
\begin{equation}\label{ho:joint-good-slice}
 \operatorname{vol}_6(\mathcal G_T^{\cap}\cap\{v=\text{constant}\})
       >(.95-.04)\mu_d(O)=.91\mu_d(O).
\end{equation}
The dependence on $v$ is measurable by the smooth finite flow maps.
The chart bounds from either witness apply at these observations
because \Cref{def:exit-test} imposes them throughout the full future
geodesic lifetime, and the matched observations are interior points
of that geodesic.

Let
\[
 I_T=\{v\in T+J_0:|\cos(\lambda v)|\ge1/2\}.
\]
The condition occupies $2/3$ of each full period.  Only the two end
periods of an interval can be incomplete, so
$|I_T|=\frac23|J_0|+O(\lambda^{-1})>\frac12|J_0|$ for large $T$.
Integration of \eqref{ho:joint-good-slice} over $I_T$ gives
\begin{equation}\label{ho:positive-intersection}
 \nu_T(\mathcal G_T^{\cap}\cap\{v\in I_T\})
                    >.455\mu_d(O)|J_0|.
\end{equation}
Apply \Cref{prop:averaged-holonomy} separately to $g_d$ and
$\widehat g_T$.  The union of the two additional exceptional sets has
measure tending to zero.  It cannot cover
\eqref{ho:positive-intersection}.  At a remaining parameter both
$|D(g_d)|$ and $|D(\widehat g_T)|$ are $o(\amp)$, uniformly, whereas
\eqref{ho:persisted-signal} says their difference exceeds $\amp$.
This is a contradiction.

We have excluded nonempty interior of $\overline{E(O,B)}$.
No closedness of $E(O,B)$ was assumed.
\end{proof}

\begin{proof}[Proof of \Cref{thm:main}]
By \Cref{prop:exit-cover}, every datum in $\Ueps$ admitting the future
extension in the statement belongs to some $E(O,B)$.
The seed basis and the positive integers $B$ are countable.
By \Cref{thm:nowhere-dense}, their union is meagre in the full relative
weighted smooth topology on $\Ueps$.  Explicitly,
\[
 \mathcal R=\bigcap_{O,B}\bigl(\Ueps\setminus\overline{E(O,B)}\bigr)
\]
is a countable intersection of relatively open dense sets and consists
of future-inextendible data.
The exit cover concerns all future exits of the full developments.
Every background observation precedes the corresponding geodesic's
first comparison runout. The exact and persisted observations lie in finite interior
regions of their full developments, where the full-lifetime chart
witnesses apply. Thus the conclusion excludes every extension specified
in \Cref{thm:main}, for this single $p_{10}$ neighborhood.
\end{proof}
